\section{Timing relation: the remaining signed inequality}
No initialized timing theorem is established. The certified endpoint signs
alone imply neither a unique zero of $Q_G$ nor an ordering of that zero and
the full-network minimum. In particular the interpolated empirical times
are not used as theorem constants.

For the actual probe, put $g_E=\nabla_\theta E$ on each probe activation cell,
and let $g_G=\nabla\Phi_I$ be the repository statistic's gradient defined
above. Almost everywhere along a compatible flow, the exact discrepancy is
\[
 \Xi(\theta)=(g_E-g_G)^TF,\qquad \dot E=Q_G+\Xi.
\]
Thus one concrete sufficient signed estimate after a reached $G$ crossing is
\begin{equation}\label{initcert:timing-missing}
 0\le Q_G\le q_+,
 \qquad \Xi\le-q_+-\delta,\qquad \delta>0.
\end{equation}
This would force the full error still to decrease on that interval. All
neurons, the ungated fixed cohort defining $Q_G$, and actual training
selectors must be retained in this comparison. At probe-gate events one
must use the actual one-sided/a.e. derivative rather than extending a
single probe-cell gradient across the event.

To localize all full-error minimizers to $[a_E,b_E]$, it suffices additionally
to certify $\dot E<0$ almost everywhere from entry to $a_E$, and
$\dot E>0$ almost everywhere from $b_E$ to the final endpoint. A reached
sign-separated $G$ crossing bracket $[a_G,b_G]$ with $b_G<a_E$ would then
prove a positive timing separation without asserting uniqueness of either
crossing. These interval-wide inequalities have not been verified.
Per the requested dependency order, no timing search was substituted for
the missing initialization-to-entry proof.
